\documentclass{article}

\usepackage{corl_2026} % CoRL 2026 submission style.
% \usepackage[final]{corl_2026}
% \usepackage[preprint]{corl_2026}

\usepackage[UTF8]{ctex}
\usepackage{amsmath}
\usepackage{amsfonts}
\usepackage{amssymb}
\usepackage{booktabs}
\usepackage{microtype}
\usepackage{xcolor}
\usepackage{url}

\renewcommand{\abstractname}{摘要}
\def\keywords#1{%
  \ifdefined\isaccepted \else
    \begin{quote}
    \textbf{关键词：} #1%
    \end{quote}
  \fi
  \ifdefined\nohyperref\else\ifdefined\hypersetup
    \hypersetup{pdfkeywords={#1}}
  \fi\fi
}

\title{DreamerVLA：融合视觉-语言-动作模型与潜在世界模型}

% 初投稿版本中，CoRL 样式会自动匿名化作者信息。
\author{
  匿名作者
}

\begin{document}
\maketitle

\begin{abstract}
具身智能研究主要沿着两条路线发展：视觉-语言-动作（Vision-Language-Action, VLA）模型和基于世界模型的强化学习方法。VLA 模型受益于大型视觉语言模型提供的强语义理解能力，但通常缺乏对环境动力学建模和长时程推理的能力。相比之下，Dreamer 等世界模型方法可以通过潜在动力学模型进行长时程规划，但其感知能力和语义理解能力往往较弱。本文探索 DreamerVLA，一种结合两类方法优势的框架：它使用冻结的 VLA 编码器嵌入多模态观测，学习紧凑的潜在动力学，并通过想象 rollout 优化 actor-critic 策略。该框架旨在将语义感知与基于模型的规划结合起来，用于具身决策。
\end{abstract}

\keywords{视觉-语言-动作模型，世界模型，机器人学习}

\section{引言}

在具身智能中，智能体需要根据视觉观测、语言指令和本体感知信号预测下一段动作序列。现有方法大体可以分为两类主流范式。第一类是视觉-语言-动作（Vision-Language-Action, VLA）范式，其中策略以图像、语言指令和机器人状态为条件，自回归地生成动作序列。第二类是以 Dreamer 为代表的世界模型范式，它学习潜在动力学，在给定动作的条件下预测未来世界状态，并使用 actor-critic 框架评估候选动作、选择最终动作序列。

然而，这两类范式各有局限。VLA 模型具有较强的语义理解能力，能够根据视觉观测和语言指令生成所需动作序列，但通常缺乏显式机制来建模动作的长期后果。尤其是，VLA 模型难以推理不同动作幅度或动作变化会如何导致不同的未来结果。虽然这一局限在单步预测中未必造成严重错误，但这意味着当前动作往往是在缺乏对后续动作选择及其长期影响充分考虑的情况下生成的，从而导致较弱的长时程决策能力。相比之下，世界模型显式建模环境动力学，因此天然支持长时程推理。然而，这类方法通常语义理解能力较弱，难以通过高层语言指令控制机器人行为。此外，纯世界模型方法在真实场景中也较难扩展，因为真实数据采集成本高昂，并且学习到的世界模型与真实环境之间的不匹配会在多步预测中累积，进而造成仿真性能与真实部署之间的显著差距。

为了结合 VLA 的语义泛化能力和世界模型的动力学感知优化能力，我们提出 DreamerVLA，一个在 Dreamer 范式下训练 VLA 的统一框架。我们的关键观察是，Dreamer 中的 actor 在功能上与 VLA 策略相似，因此可以用基于 VLA 的策略替换传统 actor，并将 VLA 纳入 Dreamer 训练框架。具体而言，在训练阶段，我们首先使用预训练 VLA 编码器将文本、图像和本体感知状态等多模态输入编码为表示当前环境的嵌入向量。随后，瓶颈模块对该表示进行压缩和精炼；该模块可以由 TokenLearner 或简单 MLP 实现。得到的潜在状态被输入世界模型以预测下一环境状态，同时 actor-critic 模块被训练用于评估状态-动作对的质量并指导策略优化。在推理阶段，智能体可以利用世界模型和 actor-critic 模块进行基于想象的决策，从而获得长时程规划和推理能力。

\section{相关工作}

\paragraph{视觉-语言-动作模型。}
视觉-语言-动作模型将视觉观测、语言指令和机器人状态映射到可执行动作，为语言条件下的机器人操作提供了统一接口。这类方法通常借助大规模视觉语言模型的语义先验，在跨任务泛化和语言理解方面表现较强。然而，许多 VLA 模型主要依赖监督模仿学习，缺乏显式的未来状态建模机制，因此在需要多步规划、失败恢复或长期后果评估的任务中仍然存在局限。

\paragraph{具身控制中的世界模型。}
世界模型通过学习环境在潜在空间中的动力学，为强化学习和规划提供可微、可想象的预测模型。Dreamer 系列方法展示了在潜在空间中进行多步想象 rollout 并优化 actor-critic 策略的有效性。与直接从像素学习策略相比，世界模型能够更好地利用长期回报信号，但其语义感知和语言条件控制能力通常弱于预训练 VLA 模型。

\paragraph{面向视觉-语言-动作模型的强化学习。}
近期工作开始超越纯模仿学习，将强化学习作为 VLA 模型的后训练机制。iRe-VLA 表明，直接对大型 VLA 策略应用在线强化学习往往不稳定且计算成本高，因此在基于强化学习的探索和监督学习之间交替，以更稳定地吸收成功轨迹。RIPT-VLA 则提出一种更简单、无 critic 的交互式后训练范式，仅使用稀疏二值成功奖励，并结合动态 rollout 采样和留一法优势估计来稳定策略更新。为了改善长时程信用分配，TGRPO 在组相对优化中引入步级和轨迹级相对优势，并依赖稠密多阶段奖励提供更细粒度的监督。从系统角度看，VLA-RL 将自回归 VLA 控制表述为轨迹级在线强化学习，并引入机器人过程奖励模型以及 critic warmup、可扩展批量解码等实现策略，使大规模在线优化更加可行。作为直接环境交互的补充，VLA-RFT 利用数据驱动世界模型作为可控模拟器，并使用从想象视觉 rollout 中计算得到的经验证稠密轨迹级奖励来微调 VLA 策略。总体来看，这些研究表明，强化学习增强的 VLA 训练正在形成三条方向：将强化学习与模仿式更新交替结合、在稀疏或成形奖励下直接进行在线后训练，以及在学习到的模拟器中进行基于世界模型的强化微调。

\section{方法}

\subsection{总体框架}

我们提出 \textbf{DreamerVLA}，一种 Dreamer 风格的训练框架，将预训练视觉-语言-动作（Vision-Language-Action, VLA）模型的语义动作先验与潜在世界模型的长时程优化能力结合起来。当前 DreamerVLA 流程使用 pi0 风格的 action-query VLA head 产生动作条件隐藏观测，在这些隐藏观测上训练 DreamerV3 recurrent state-space model（RSSM），并通过想象潜在 rollout 优化 actor-critic 策略。

在每个时间步 $t$，机器人观测两路 RGB 图像、语言指令和本体感知状态：
\[
o_t = \{I_t^{\text{third}}, I_t^{\text{wrist}}, x_t^{\text{prop}}, l\}.
\]
VLA 模型首先将该多模态观测转换为 action-query hidden 序列
\[
u_t = f_{\text{VLA}}^{\text{query}}(o_t) \in \mathbb{R}^{K \times 1024},
\]
其中 $K$ 是动作预测长度。随后将该序列展平为世界模型观测：
\[
e_t = \mathrm{vec}(u_t) \in \mathbb{R}^{5120},
\]
在 LIBERO 配置中我们使用 $K=5$。这些隐藏观测被离线预计算，并以 sidecar HDF5 文件的形式与原始 demonstrations 对齐保存，因此世界模型和 actor-critic 训练不会反向传播到大型 VLA 主干。

DreamerVLA 包含三个可训练组件：
\begin{enumerate}
    \item \textbf{action-hidden DreamerV3 世界模型}：使用 pi0 action-query hidden state 作为离散 RSSM 的观测，并同时重构像素和 action-hidden observation；
    \item \textbf{pi0 action-hidden actor}：将想象 RSSM feature 映射回 action-query hidden state，并复用预训练 pi0 output projection 产生动作块；
    \item \textbf{two-hot critic}：由 DreamerV3 风格的想象 $\lambda$-return 进行训练。
\end{enumerate}

该设计将 VLA 主干固定为语义动作条件编码器，在紧凑潜在空间中学习任务动力学，并在无需额外环境交互的情况下通过想象 rollout 优化 action head。

\subsection{整体训练流程}

DreamerVLA 采用四阶段训练流程。第一阶段，我们对 LIBERO demonstrations 进行数据准备，去除 no-op transition，并保存同步的第三视角 RGB、腕部 RGB、机器人状态、动作、奖励以及终止标记。动作预测长度设为 $K=5$，每个动作是 7 维连续控制向量。

第二阶段，我们通过监督动作预测对 VLA policy 进行 warm-start。在预训练 RynnVLA/Chameleon 主干的基础上，我们在 LIBERO demonstrations 上微调 pi0 风格的 action-query head。该 head 不再直接从全部语言-图像 token 中预测动作，而是为未来每个动作步追加一个可学习 action query，并从这些 query 的隐藏状态中回归一个动作 chunk。该阶段得到的监督 VLA checkpoint 会作为后续 DreamerVLA 组件的初始化。

第三阶段，我们冻结 VLA 主干，并为每一帧离线预计算 action-conditioned hidden observation。生成的 sidecar 文件存储 $e_t$，同时记录 VLA checkpoint、action horizon、action-head type 和 hidden-source type 等元数据。数据加载器会在启动时检查这些元数据，从而避免将过期 sidecar 与不兼容的 checkpoint 混用。

第四阶段，我们训练 action-hidden DreamerV3 世界模型，并进一步通过想象 rollout 优化 actor 和 critic。世界模型预训练可以作为独立阶段运行，其 checkpoint 随后加载到 joint DreamerVLA 训练中。在 actor-critic 阶段，replay window 中推断出的 posterior latent state 作为想象轨迹起点；actor 从想象隐藏状态产生动作 chunk，并使用每个 chunk 的第一步动作推进 RSSM rollout。

\subsection{冻结的 VLA Action-Query 特征}

我们使用预训练 RynnVLA/Chameleon 模型作为冻结的多模态编码器。对于每一帧，第三视角图像、腕部相机图像、任务文本和可选本体感知状态被序列化为 VLA 模型使用的 prompt 格式。Chameleon 图像 tokenizer 和文本 tokenizer 将输入转换为统一 token 序列，并由冻结 Transformer 主干处理。

pi0 风格的 action-query head 会追加 $K$ 个可学习 query token，每个 query 对应一个未来动作步。上下文 token 通过 attention 条件化这些 query，head 将 query 输出作为紧凑的动作条件隐藏状态：
\[
u_t = [u_t^1,\ldots,u_t^K],
\qquad
u_t^k \in \mathbb{R}^{1024}.
\]
世界模型使用其展平表示 $e_t=\mathrm{vec}(u_t)$ 作为观测。相比存储完整 VLA token sequence，该表示更小，并且与下游 action head 直接对齐。Dreamer 训练过程中，VLA 主干始终冻结，不进入优化计算图。

\subsection{Action-Hidden DreamerV3 世界模型}

世界模型遵循 DreamerV3 recurrent state-space model，但将标准像素编码器替换为作用于预计算 action-hidden observation 的轻量编码器。给定 action-hidden observation 序列 $\{e_t\}_{t=1}^{T}$ 和动作序列 $\{a_t\}_{t=1}^{T}$，编码器首先将每个 $e_t$ 映射到 RSSM 观测空间。RSSM 状态定义为
\[
s_t = (h_t, z_t),
\]
其中 $h_t \in \mathbb{R}^{8192}$ 是确定性循环状态，$z_t$ 是由 $32$ 个随机变量组成、每个变量具有 $64$ 个类别的离散随机状态。展平后的潜在特征为
\[
\phi_t = [h_t; \mathrm{vec}(z_t)] \in \mathbb{R}^{10240}.
\]

\paragraph{RSSM 转移与后验。}
确定性转移是 DreamerV3 block-GRU 更新，以前一确定性状态、随机状态和动作为条件：
\[
h_t = g_\theta(h_{t-1}, z_{t-1}, a_t).
\]
模型预测先验分布 $p_\theta(z_t \mid h_t)$，并根据观测得到后验：
\[
q_\theta(z_t \mid h_t, e_t).
\]
我们采用 DreamerV3 的 unimix 类别参数化，并将 KL 正则项拆分为 dynamics 和 representation 两部分：
\[
\mathcal{L}_{\text{dyn}} =
\mathrm{KL}\big(\mathrm{sg}(q_\theta(z_t \mid h_t,e_t)) \,\|\, p_\theta(z_t \mid h_t)\big),
\]
\[
\mathcal{L}_{\text{rep}} =
\mathrm{KL}\big(q_\theta(z_t \mid h_t,e_t) \,\|\, \mathrm{sg}(p_\theta(z_t \mid h_t))\big).
\]

\paragraph{多目标重构。}
RSSM 同时使用视觉重构和 action-hidden 重构进行训练。像素解码器从 $\phi_t$ 重构拼接后的第三视角与腕部图像：
\[
\hat{I}_t = d_{\text{img}}(\phi_t),
\qquad
\mathcal{L}_{\text{img}} = \|\hat{I}_t - I_t\|_2^2.
\]
隐藏状态解码器预测展平后的 pi0 action-query hidden observation：
\[
\hat{e}_t = d_{\text{hid}}(\phi_t),
\qquad
\mathcal{L}_{\text{hid}} = \|\hat{e}_t - e_t\|_2^2.
\]
该隐藏重构连接了 RSSM 与 pi0 action-hidden actor：在想象阶段，actor 消费 $\hat{e}_t$，而不是原始像素或完整 VLA token state。

\paragraph{奖励与 continuation。}
世界模型从潜在特征预测稀疏任务奖励和 continuation：
\[
\hat{r}_t = d_r(\phi_t), \qquad \hat{c}_t = d_c(\phi_t).
\]
对于 LIBERO 风格的稀疏成功奖励，当前实现使用 binary reward head；continuation head 则通过非终止目标上的二元交叉熵进行训练。

完整世界模型目标为
\[
\mathcal{L}_{\text{WM}}
=
\lambda_{\text{img}}\mathcal{L}_{\text{img}}
+ \lambda_{\text{dyn}}\mathcal{L}_{\text{dyn}}
+ \lambda_{\text{rep}}\mathcal{L}_{\text{rep}}
+ \lambda_{\text{hid}}\mathcal{L}_{\text{hid}}
+ \lambda_r\mathcal{L}_{r}
+ \lambda_c\mathcal{L}_{c}.
\]

\subsection{Pi0 Action-Hidden Actor}

DreamerVLA 不使用小型 MLP 策略，而是复用监督训练得到的 pi0 风格 VLA head 中的动作先验。世界模型首先将想象潜在特征解码为重构 action-hidden 向量：
\[
\hat{e}_t = d_{\text{hid}}(\phi_t) \in \mathbb{R}^{5120},
\]
随后 reshape 为
\[
\hat{u}_t \in \mathbb{R}^{K \times 1024}.
\]
一个轻量 adapter 对该序列进行修正，预训练 pi0 output projection 则将每个 action-query hidden state 映射为连续动作。因此 actor 输出动作块：
\[
\mu_{t:t+K-1} = \pi_\psi(\hat{u}_t) \in \mathbb{R}^{K \times A}.
\]
在 LIBERO 配置中，$K=5$，$A=7$。想象训练时，我们取动作块中的第一步作为动作均值：
\[
\mu_t = \mu_{t:t+K-1}[0],
\]
并使用可学习的对角 log 标准差定义高斯策略：
\[
a_t \sim \mathcal{N}(\mu_t, \mathrm{diag}(\sigma^2)).
\]
actor 可以选择冻结或微调 pi0 output projection。两种情况下，策略都由监督 VLA checkpoint 初始化，并在之后通过 DreamerV3 actor-critic 更新继续优化。

\subsection{基于想象的 Actor-Critic 学习}

Actor-critic 训练从 replay 序列推断得到的后验潜在状态开始。从每个起始状态出发，世界模型生成长度为 $H$ 的想象轨迹：
\[
a_t \sim \pi_\psi(\cdot \mid \phi_t),
\qquad
s_{t+1} \sim p_\theta(\cdot \mid s_t,a_t),
\qquad
\hat{r}_{t+1}=d_r(\phi_{t+1}), \quad
\hat{c}_{t+1}=d_c(\phi_{t+1}).
\]
梯度在想象世界模型状态上停止，这与标准 DreamerV3 actor 更新一致。

critic 接收 RSSM 特征 $\phi_t$，并预测 symlog value bins 上的类别分布。其标量价值为 symexp 变换后的期望：
\[
V_\omega(\phi_t)=\mathrm{symexp}\left(\mathbb{E}_{b\sim p_\omega(\cdot\mid \phi_t)}[b]\right).
\]
Polyak 平均的目标 critic 提供 bootstrap value。我们计算 continuation 加权的 $\lambda$-return：
\[
R_t^\lambda =
\hat{r}_{t+1}
+ \gamma \hat{c}_{t+1}
\left((1-\lambda)V_{\bar{\omega}}(\phi_{t+1})
+ \lambda R_{t+1}^{\lambda}\right).
\]

actor 使用基于想象 advantage 的 likelihood-ratio 目标进行训练：
\[
\mathcal{L}_{\pi}
=
-
\mathbb{E}\left[
w_t
\left(
\log \pi_\psi(a_t\mid \phi_t)\,
\mathrm{sg}\left(\frac{R_t^\lambda - V_\omega(\phi_t)}{S}\right)
+ \eta \mathcal{H}[\pi_\psi(\cdot\mid \phi_t)]
\right)
\right],
\]
其中 $w_t$ 是累计 continuation discount，$S=\max(1,P_{95}^{\text{ema}}-P_{5}^{\text{ema}})$ 是运行百分位回报尺度。critic 则最大化 two-hot symlog 目标的对数概率：
\[
\mathcal{L}_{V}
=
-
\mathbb{E}\left[
\log p_\omega\left(\mathrm{twohot}(\mathrm{sg}(R_t^\lambda)) \mid \phi_t\right)
\right].
\]

\subsection{训练策略}

训练在同一批 replay windows 上交替执行两个阶段。首先，世界模型使用真实图像序列、预计算 action-hidden observation、动作、奖励和终止标签进行更新。随后，actor 和 critic 使用学习到的 RSSM 生成的想象 rollout 进行更新。整个过程中，只有世界模型、pi0 action-hidden actor 和 critic 被优化；用于生成 action-query hidden observation 的 VLA 主干始终保持冻结。

\section{实验}

\subsection{实验设置}

\paragraph{基线方法。}

\subsection{主要结果}

\begin{table}[t]
\centering
\caption{在 LIBERO 上与代表性方法的比较。}
\label{tab:libero_main_simple_zh}
\setlength{\tabcolsep}{6pt}
\begin{tabular}{lccccc}
\toprule
方法 & Spatial & Object & Goal & Long & Avg. \\
\midrule
GR-1              & -- & -- & -- & -- & -- \\
RT-1              & -- & -- & -- & -- & -- \\
Diffusion Policy  & -- & -- & -- & -- & -- \\
Octo              & -- & -- & -- & -- & -- \\
OpenVLA           & -- & -- & -- & -- & -- \\
RynnVLA           & -- & -- & -- & -- & -- \\
\midrule
\textbf{DreamerVLA（本文方法）} & -- & -- & -- & -- & -- \\
\bottomrule
\end{tabular}
\end{table}

\subsection{消融实验}

\begin{table}[t]
\centering
\caption{LIBERO 上的消融实验。每一行在基线基础上加入一个组件。}
\label{tab:libero_ablation_zh}
\setlength{\tabcolsep}{5pt}
\begin{tabular}{lcccccccc}
\toprule
变体 & WM & PPO & Adv. & KL & Spatial & Goal & Long & Avg. \\
\midrule
Baseline（仅 VLA）       &   &   &   &   & -- & -- & -- & -- \\
+ PPO                    &   & \checkmark &   &   & -- & -- & -- & -- \\
+ PPO + WM               & \checkmark & \checkmark &   &   & -- & -- & -- & -- \\
+ PPO + WM + Adv.        & \checkmark & \checkmark & \checkmark &   & -- & -- & -- & -- \\
完整模型                 & \checkmark & \checkmark & \checkmark & \checkmark & -- & -- & -- & -- \\
\bottomrule
\end{tabular}
\end{table}

\subsection{泛化性与鲁棒性}

\section{结论}

本文提出 DreamerVLA，一个结合预训练视觉-语言-动作模型语义表示能力与潜在世界模型长时程推理能力的框架。通过冻结 VLA 编码器，并训练紧凑的潜在动力学模型以及基于想象的 actor-critic 模块，DreamerVLA 将语言锚定的感知与任务相关的预测控制分离开来。该设计避免了像素级重构，同时支持通过想象潜在 rollout 进行策略优化。未来工作将进一步在更广泛的真实机器人操作任务上评估该框架，并研究如何利用世界模型不确定性提升安全部署能力。

\clearpage

% 中文版本目前没有真实引用，待补充 bib 条目和 \cite{} 后再启用参考文献。
% CoRL 样式会自动使用 corlabbrvnat，无需手动设置 \bibliographystyle。
% \bibliography{example}

\end{document}
